\section{Dialect-Trap Taxonomy and Corpus}\label{sec:taxonomy}

\textit{Taxonomy.} We derived the taxonomy from the ASE and SQL Server language
references~\cite{aseref,mstsql} and from the authors' migration experience. Each category is labelled by
\emph{observability}: how an unconverted or naively converted construct manifests on SQL Server.
\emph{Compile-visible} traps (C) fail at deployment; \emph{silent} traps (S) deploy and run but change
behaviour; \emph{deprecation} traps (D) behave identically but use features scheduled for removal; and
\emph{control} cases (K) contain already-portable code and measure harmful over-editing.
Table~\ref{tab:taxonomy} summarises the categories and corpus.

\maybeinput{tables/corpus.tex}

\textit{Corpus.} Each case contains (i) the ASE source object a converter sees, (ii) a SQL Server setup
script with target-side tables and data, (iii) a reviewed reference translation, and (iv) one or more calls
whose outcomes define behaviour, chosen to exercise boundary inputs such as NULL arguments, empty strings,
missing keys, and month-end dates. The development corpus holds \NumDevCases{} cases, of which \NumSilentCases{} are silent and
\NumControlCases{} are controls; each isolates one trap and is short. A further \NumHeldout{}
\emph{held-out} composite procedures (7--22 lines, two to five trap categories each, such as a funds transfer or a
batched purge) were written after the trap checklist was frozen and are reported separately. All cases are
synthetic and written by the authors. Two cases encode a migration-policy decision rather than a single
correct behaviour (the result shape that replaces \texttt{COMPUTE BY}, and whether to preserve ASE's
empty-string-as-blank rule); these \NumPolicyCases{} are scored on G1/G2 only, and every other case is scored
differentially on at least two inputs.

Fig.~\ref{fig:trap} shows a silent trap. The ASE procedure keeps customers without large orders; a
translation that moves \texttt{o.amt > 30} into \texttt{WHERE} compiles, passes lint, and drops them.

\begin{figure}[t]
\begin{lstlisting}
-- ASE source
select c.id, c.name, o.amt from cust c, ord o
where c.id *= o.cust_id and o.amt > 30
-- Plausible but wrong (inner-join semantics)
SELECT c.id, c.name, o.amt FROM cust c
LEFT JOIN ord o ON c.id = o.cust_id WHERE o.amt > 30
-- Reference: restriction stays in the join
SELECT c.id, c.name, o.amt FROM cust c
LEFT JOIN ord o ON c.id = o.cust_id AND o.amt > 30
\end{lstlisting}
\caption{A silent outer-join trap (corpus case C02).}
\label{fig:trap}
\end{figure}

\textit{Reference validation.} Every reference translation must deploy on SQL Server 2022 and produce its
expected outcomes before any run starts; a failing reference aborts the experiment. For the
\NumReviewCases{} cases whose expected behaviour depends on ASE-specific semantics, that behaviour is
derived from the ASE 16 documentation~\cite{aseref} under default server options, and the cases are flagged
in the artifact for re-checking against a live instance \note{if you execute them on ASE, say so here with
the version, e.g. ``and confirmed by executing the sources on ASE 16.0 SP04''}.
